\section{Análisis horizontal}

El análisis horizontal compara una misma partida contable entre dos o más periodos. Su finalidad es identificar aumentos, disminuciones y tendencias.

Las fórmulas son:

\[
\text{Variación absoluta} = \text{Valor actual} - \text{Valor base}
\]

\[
\text{Variación relativa} =
\frac{\text{Variación absoluta}}{\text{Valor base}} \times 100
\]

El procedimiento consiste en seleccionar la partida, elegir un periodo base, calcular la variación absoluta, calcular la variación relativa e interpretar el resultado.

Este método responde a la pregunta: \textbf{¿cuánto cambió una partida entre periodos?}

\subsection{Ejemplo básico de análisis horizontal}

Enunciado: una empresa desea comparar sus ventas entre dos periodos para medir el crecimiento comercial.

Datos:

\begin{itemize}
  \item Ventas 2024: S/ 100,000.00.
  \item Ventas 2025: S/ 120,000.00.
\end{itemize}

Cálculo:

\[
\text{S/ 120,000.00} - \text{S/ 100,000.00} = \text{S/ 20,000.00}
\]

La variación porcentual es:

\[
\frac{\text{S/ 20,000.00}}{\text{S/ 100,000.00}} \times 100 = 20.00\%
\]

Interpretación: las ventas aumentaron S/ 20,000.00, equivalentes al 20.00\% del periodo base.

\subsection{Ejemplo aplicado de análisis horizontal}

Enunciado: una empresa de servicios digitales contrata infraestructura cloud para alojar su plataforma. Al comparar dos periodos, la gerencia observa un incremento importante del costo mensual.

Datos:

\begin{itemize}
  \item Costo cloud 2024: S/ 8,000.00.
  \item Costo cloud 2025: S/ 12,000.00.
\end{itemize}

Cálculo:

\[
\text{S/ 12,000.00} - \text{S/ 8,000.00} = \text{S/ 4,000.00}
\]

\[
\frac{\text{S/ 4,000.00}}{\text{S/ 8,000.00}} \times 100 = 50.00\%
\]

Interpretación y decisión: el incremento puede estar justificado si creció la demanda de usuarios; sin embargo, si los ingresos no crecieron en la misma proporción, debe revisarse la eficiencia de la arquitectura, proveedores o plan contratado.
